\documentclass[11pt]{article}

% Change "review" to "final" to generate the final (sometimes called camera-ready) version.
% Change to "preprint" to generate a non-anonymous version with page numbers.
\usepackage[review]{acl}
% \usepackage{acl}

% Standard package includes
\usepackage{times}
\usepackage{latexsym}

\usepackage{booktabs}

% For proper rendering and hyphenation of words containing Latin characters (including in bib files)
% \usepackage[T1]{fontenc}
\usepackage[english]{babel}
%\usepackage[T1,T2A]{fontenc}
\usepackage[T2A, T1]{fontenc}
\usepackage{inconsolata}
% For Vietnamese characters
% \usepackage[T5]{fontenc}
% See https://www.latex-project.org/help/documentation/encguide.pdf for other character sets

% This assumes your files are encoded as UTF8
\usepackage[utf8]{inputenc}

% This is not strictly necessary, and may be commented out,
% but it will improve the layout of the manuscript,
% and will typically save some space.
\usepackage{microtype}

% This is also not strictly necessary, and may be commented out.
% However, it will improve the aesthetics of text in
% the typewriter font.
\usepackage{inconsolata}

%Including images in your LaTeX document requires adding
%additional package(s)
\usepackage{graphicx}

\usepackage{amsmath}
\usepackage{alltt}

% Helper for short Ukrainian/Cyrillic inline snippets without leaking font settings.
\newcommand{\uk}[1]{{\fontencoding{T2A}\selectfont #1}}

% If the title and author information does not fit in the area allocated, uncomment the following
%
%\setlength\titlebox{<dim>}
%
% and set <dim> to something 5cm or larger.

\title{Bridging the Skill Gap: An LLM Pipeline and Large-Scale Dataset for ESCO Skill Mapping
}

\author{
  \textbf{Ihor Pysmennyi\textsuperscript{1}},
  \textbf{Roman Kyslyi\textsuperscript{2}},
  \textbf{Taras Baraniuk\textsuperscript{2}}
\\
\textsuperscript{1}National Technical University of Ukraine "Igor Sikorsky Kyiv Polytechnic Institute" \\
\textsuperscript{2}Kyiv School of Economics
\\
\small{
  \textbf{Correspondence:}   \href{mailto:pysmennyi.ihor@lll.kpi.ua}{pysmennyi.ihor@lll.kpi.ua}
  \href{mailto:rkyslyi@kse.org.ua}{rkyslyi@kse.org.ua}, 
  \href{mailto:tbaraniuk@kse.org.ua}{tbaraniuk@kse.org.ua}, 
}
}

% Author information can be set in various styles:
% For several authors from the same institution:
% \author{Author 1 \and ... \and Author n \\
%         Address line \\ ... \\ Address line}
% if the names do not fit well on one line use
%         Author 1 \\ {\bf Author 2} \\ ... \\ {\bf Author n} \\
% For authors from different institutions:
% \author{Author 1 \\ Address line \\  ... \\ Address line
%         \And  ... \And
%         Author n \\ Address line \\ ... \\ Address line}
% To start a separate ``row'' of authors use \AND, as in
% \author{Author 1 \\ Address line \\  ... \\ Address line
%         \AND
%         Author 2 \\ Address line \\ ... \\ Address line \And
%         Author 3 \\ Address line \\ ... \\ Address line}

% \author{First Author \\
%   Affiliation / Address line 1 \\
%   Affiliation / Address line 2 \\
%   Affiliation / Address line 3 \\
%   \texttt{email@domain} \\\And
%   Second Author \\
%   Affiliation / Address line 1 \\
%   Affiliation / Address line 2 \\
%   Affiliation / Address line 3 \\
%   \texttt{email@domain} \\}

%\author{
%  \textbf{First Author\textsuperscript{1}},
%  \textbf{Second Author\textsuperscript{1,2}},
%  \textbf{Third T. Author\textsuperscript{1}},
%  \textbf{Fourth Author\textsuperscript{1}},
%\\
%  \textbf{Fifth Author\textsuperscript{1,2}},
%  \textbf{Sixth Author\textsuperscript{1}},
%  \textbf{Seventh Author\textsuperscript{1}},
%  \textbf{Eighth Author \textsuperscript{1,2,3,4}},
%\\
%  \textbf{Ninth Author\textsuperscript{1}},
%  \textbf{Tenth Author\textsuperscript{1}},
%  \textbf{Eleventh E. Author\textsuperscript{1,2,3,4,5}},
%  \textbf{Twelfth Author\textsuperscript{1}},
%\\
%  \textbf{Thirteenth Author\textsuperscript{3}},
%  \textbf{Fourteenth F. Author\textsuperscript{2,4}},
%  \textbf{Fifteenth Author\textsuperscript{1}},
%  \textbf{Sixteenth Author\textsuperscript{1}},
%\\
%  \textbf{Seventeenth S. Author\textsuperscript{4,5}},
%  \textbf{Eighteenth Author\textsuperscript{3,4}},
%  \textbf{Nineteenth N. Author\textsuperscript{2,5}},
%  \textbf{Twentieth Author\textsuperscript{1}}
%\\
%\\
%  \textsuperscript{1}Affiliation 1,
%  \textsuperscript{2}Affiliation 2,
%  \textsuperscript{3}Affiliation 3,
%  \textsuperscript{4}Affiliation 4,
%  \textsuperscript{5}Affiliation 5
%\\
%  \small{
%    \textbf{Correspondence:} \href{mailto:email@domain}{email@domain}
%  }
%}

\begin{document}
\maketitle
\begin{abstract}
We present a retrieval-augmented LLM pipeline for extracting skills from labor market documents and normalizing them to the European Skills, Competences, Qualifications and Occupations (ESCO) taxonomy. The pipeline operates in three stages: (1) LLM-based skill extraction, (2) retrieval-augmented global mapping with ESCO graph enrichment, and (3) context-aware per-document rescoring. By design, the architecture is language-agnostic: it requires only an ESCO locale file and target-language prompts, making it applicable to any of the 28 languages in which ESCO is available. We instantiate and evaluate the pipeline on Ukrainian, a morphologically rich, lower-resource language, and apply it to [N] job vacancies and [N] resumes collected from Work.ua, the largest Ukrainian job platform. We release two complementary resources: (1) a raw corpus of [~200K] Ukrainian labour market documents and (2)~an enriched version of the same corpus in which each document is augmented with structured ESCO skill annotations, including mapped URIs, labels, confidence scores, and extraction metadata. We further demonstrate the dataset’s utility through a preliminary skill supply–demand gap analysis. The pipeline code, datasets, and evaluation materials are publicly available.
\footnote{{GitHub and HuggingFace URLs}}

\end{abstract}

\section{Introduction}

Understanding the dynamics of labour markets--what skills employers demand, what skills workers offer, and where the gaps lie--is a prerequisite for evidence-based workforce policy, educational planning, and career guidance. Online job platforms generate large volumes of data that could, in principle, support such analyses at scale. However, raw skill mentions in job postings and resumes are heterogeneous: they combine free-text descriptions, platform-specific tags, abbreviations, and domain jargon, all expressed in the morphological forms of the source language. To enable cross-document, cross-platform, and cross-temporal comparison, these mentions must be normalized to a controlled vocabulary.

The European Skills, Competences, Qualifications and Occupations (ESCO) taxonomy provides such a vocabulary. Maintained by the European Commission, ESCO v1.2.1 covers approximately 13,890 skill and competence concepts organized as a directed acyclic graph, with preferred and alternative labels in 28 languages. Since 2022, ESCO has included full Ukrainian localization, a development motivated by the need to support labour-market integration for displaced Ukrainians across Europe \citep{EuroComission2022}. ESCO’s multilingual coverage and hierarchical structure make it an attractive normalization target, but the sheer scale of the taxonomy, nearly 14,000 leaf-level skills, makes direct classification intractable for standard supervised methods and exceeds the context windows of most language models.
We address this challenge with a three-stage pipeline that decomposes the problem into manageable steps. In Stage 1, a locally hosted LLM extracts skill mentions from each document. In Stage 2, a retrieval-augmented mapping module constructs a ranked candidate list of ESCO skills using fuzzy matching, embedding similarity, and ESCO graph traversal, then prompts the LLM to select the best match for each skill from this narrowed candidate set. In Stage 3, a context-aware rescoring pass re-evaluates low-confidence and unmapped skills against the specific document context, resolving ambiguities that the context-free global mapping could not. Crucially, the architecture is designed to be language-agnostic: given an ESCO locale and translated prompts, the pipeline can be applied to any ESCO-supported language without retraining.
We instantiate this pipeline for Ukrainian—a morphologically rich, relatively lower-resource language that presents particular challenges for skill extraction due to extensive inflection, prevalent transliteration of technical terms, and limited availability of domain-specific NLP tools. We apply it to a corpus of approximately100K job vacancies and 100K resumes collected from Work.ua, the largest Ukrainian employment platform. The result is two complementary datasets released on HuggingFace: a raw corpus of Ukrainian labour-market documents, and ESCO-aligned skill annotations derived from those documents by our pipeline.

Our contributions are as follows:

\begin{itemize}

\item \textbf{A reusable, language-agnostic pipeline} for ESCO skill extraction and normalization, combining retrieval-augmented LLM mapping with ESCO graph enrichment and context-aware rescoring. The pipeline operates on a locally hosted open-weights LLM and pre-computed embeddings, requiring no dependence on proprietary APIs at inference time.

\item \textbf{Two large-scale Ukrainian datasets:} (a)~a raw corpus of {\raise.17ex\hbox{$\scriptstyle\sim$}}200K job vacancies and resumes from Work.ua, and (b)~an analysis-ready enriched corpus where each document is augmented with an array of ESCO skill annotations (extracted skill string, mapped URI, label, confidence, method). Researchers can immediately compute skill demand by region, skill supply by experience level, or vacancy--resume skill overlap without running any extraction pipeline. To our knowledge, this is the first publicly released Ukrainian labour-market corpus with standardized skill annotations.

\item \textbf{A multi-faceted evaluation} that addresses the impracticality of exact-URI gold-standard annotation at ESCO's scale, instead combining lightweight expert review of high-frequency mappings, cross-document consistency analysis, intrinsic quality metrics, ablation across pipeline stages, and comparison with simpler baselines.

\item \textbf{A preliminary skill gap analysis} demonstrating the dataset's utility for identifying mismatches between skill supply (resumes) and demand (vacancies) in the Ukrainian labour market.

\end{itemize}
The remainder of this paper is organized as follows. ADD LATER


\section{Related work}
\label{sec:related}


Our work draws on two lines of research: skill extraction and ESCO-based normalization from job descriptions, and ESCO-driven labour-market intelligence with particular attention to Ukrainian data.

\subsection{Skill Extraction and ESCO Mapping}
\label{sec:related-mapping}

Mapping unstructured skill mentions to ESCO's 13,890-node taxonomy is an active and rapidly evolving problem. We trace the field's progression from embedding-based matching through contrastive training to LLM-augmented pipelines, highlighting how each advance addresses a limitation of its predecessors--and where open challenges remain.

\paragraph{Embedding-based matching.}
The simplest approach to ESCO normalization is direct vector similarity: embed the extracted skill and every ESCO label, then select the nearest neighbour. The ESCOX library \citep{Kavargyris2025} implements this with transformer embeddings and cosine similarity, and Nesta's Skills Extractor Library \citep{nesta} pairs span extraction with a similar taxonomy lookup. These systems are fast and require no LLM at inference time, but they treat ESCO as a flat list of targets--ignoring the taxonomy's hierarchical structure--and struggle when surface forms diverge from official labels (e.g., Ukrainian domain jargon or transliterated terms). Our \texttt{EmbeddingMapper} baseline is conceptually equivalent.

\paragraph{Contrastive training and extreme multi-label classification.}
\citet{Decorte2023} reframe the problem as extreme multi-label classification and use LLMs to generate 138K synthetic (skill, sentence) pairs covering 99.5\% of ESCO skills. Contrastive training on this synthetic data yields 15--25 point improvements in R-Precision@5 over embedding-only baselines, evaluated on SkillSpan-ESCO \citep{Zhang2022} -- the standard English benchmark for skill extraction, which provides 14,500 expert-annotated sentences. Building on this, \citet{Decorte2025} propose ConTeXT-match, a contrastive bi-encoder with token-level attention that achieves state-of-the-art results while being substantially faster than LLM-based approaches. They also release Skill-XL, a new benchmark with exhaustive sentence-level annotations. Separately, \citet{zhang2022skillextractionjobpostings} show that ESCO labels can serve directly as weak supervision by matching embedded n-grams to embedded ESCO skills, avoiding manual annotation entirely. This line of work demonstrates a clear trend toward encoder-efficient architectures optimized for real-time inference. However, the contrastive models require large-scale synthetic training data generated for each target language, and like embedding-based methods, they operate over a flat ESCO target space without exploiting the taxonomy's graph structure.

\paragraph{LLM-augmented retrieve-and-rerank.}
Recent systems converge on a \emph{retrieve-then-re-rank} paradigm: narrow the search space via retrieval, then apply an LLM to select the best match from a manageable candidate set. \citet{Clavi2023} generate synthetic data to train a contrastive retriever, then use GPT-4 to re-rank the top candidates -- achieving strong zero-shot performance without manual annotations. However, their pipeline processes each document independently: retrieval and re-ranking cost scales linearly with corpus size. SkiLLMo \citep{Malandri2025} addresses the extraction bottleneck by fine-tuning a BERT model to filter skill-dense sentences before LLM extraction, reporting 91\% extraction precision and 79\% end-to-end precision. Their normalization step, however, still relies on a standard vector store. SkillGPT \citep{li2023skillgptrestfulapiservice} uses a LLaMA-based model for summarization paired with vector similarity, arguing that direct LLM prompting for ESCO labels is too slow for production use. Across all three systems, two architectural limitations persist:

\begin{enumerate}
\setlength\itemsep{0.2em}
\item \textbf{Flat candidate retrieval.} Candidates are retrieved from ESCO as an unstructured set of labels. None exploit the taxonomy's DAG structure, missing semantically adjacent concepts (parents, siblings) that may better match a raw mention than any directly retrieved label.

\item \textbf{Per-document processing.} Retrieval and LLM mapping are repeated for every document occurrence of a skill. When common skills like \uk{Продаж} (Sales) or \uk{Вміння працювати в команді} (teamwork) appear across thousands of documents, this creates significant computational redundancy.
\end{enumerate}

\paragraph{Our approach.}
Our pipeline shares the retrieve-then-re-rank paradigm but makes different trade-offs suited to offline corpus annotation at scale. \textbf{Extraction} uses prompt-based LLM generation requiring no task-specific fine-tuning or annotated training data -- only a target-language prompt -- which simplifies adaptation to new ESCO locales. \textbf{Candidate generation} combines fuzzy and embedding retrieval with ESCO graph enrichment: parent and sibling nodes are traversed via the DAG and added to the candidate pool, expanding recall for skills whose surface forms diverge from official labels. \textbf{Global deduplication} amortizes the expensive retrieval and LLM mapping to one call per unique skill string across the full corpus, rather than one per document occurrence. \textbf{Context-aware rescoring} (Stage~3) then re-introduces document-level context selectively, only for low-confidence or unmapped skills, recovering the accuracy that global deduplication would otherwise sacrifice for ambiguous terms. Our setting also differs from the encoder-efficiency trend \citep{Decorte2025}: we prioritize mapping quality over inference speed, since our pipeline runs as an offline batch process for dataset creation rather than a real-time service. Table~\ref{tab:comparison} summarizes these architectural differences.

\begin{table*}[t]
\centering
\small
\begin{tabular}{@{}lccc@{}}
\toprule
& \textbf{Extraction} & \textbf{Candidate retrieval} & \textbf{Disambiguation} \\
\midrule
\citet{Kavargyris2025} &  --  & Flat cosine similarity  & None \\
 \citet{nesta}&  &   &  \\[4pt]

\citet{Decorte2023, Decorte2025} & Trained classifier & Trained semantic retrieval  & None \\[4pt]
\citet{Clavi2023} & Trained classifier & Trained semantic retrieval  & LLM re-rank (no context) \\[4pt]
\citet{Malandri2025} & BERT filter + LLM & Flat cosine similarity & None \\[4pt]
\textbf{Ours} & Zero-shot LLM & Fuzzy + embedding  & \textbf{Context-aware} \\
& (no fine-tuning) & \textbf{+ ESCO graph}  & \textbf{rescoring (Stage~3)} \\
\bottomrule
\end{tabular}
\caption{Architectural comparison of ESCO skill mapping systems. Bold indicates novel elements introduced in our pipeline.}
\label{tab:comparison}
\end{table*}

\subsection{Ukrainian ESCO Localization and Labour-Market Analytics}
\label{sec:related-ukraine}

The European Commission released ESCO with Ukrainian localization in 2022, explicitly to facilitate skills assessment and labour-market integration for Ukrainians \citep{EuroComission2022}. This provides the preferred and alternative labels we use for fuzzy matching and embedding computation, and enables our pipeline to operate entirely in Ukrainian without cross-lingual transfer.

The closest existing large-scale effort to map Ukrainian job data to ESCO is the European Training Foundation's (ETF) Big Data for Labour Market Intelligence project \citep{Sarioglo2020}. Over eight months, they collected hundreds of thousands of online job vacancies in Ukraine, translated ESCO skill labels, and linked non-ESCO skills to ESCO entries. They note explicitly that online job vacancies contain emerging skills lacking direct ESCO counterparts. However, neither the data nor the models from this project are publicly released; the results are available only through aggregated dashboards. Our work fills this reproducibility gap by releasing both the raw corpus and an enriched version with document-level ESCO skill annotations, along with the open-source pipeline.

At the European level, the Skills~OVATE dashboard and Europass Job and Skill Trends Tool use ESCO-classified vacancy data to monitor skill demand across countries \citep{skills-ovate}.  These initiatives demonstrate the policy value of ESCO-normalized labour-market data but operate at aggregate European scales. Our Ukrainian-focused corpus complements these efforts by providing document-level, openly accessible data for a specific national labour market -- one that is underrepresented in existing open resources despite its policy significance.

\subsection{Positioning}
\label{sec:related-positioning}

Our work is positioned at the intersection of these two lines of research, with a focus on producing a reusable architecture and an open data resource rather than advancing state-of-the-art metrics on existing benchmarks.

Architecturally, we build on the retrieve-then-re-rank paradigm established by \citet{Clavi2023} and the multi-stage design of SkiLLMo \citep{Malandri2025}, but our design goals differ (see Table~\ref{tab:comparison}). Where prior systems optimize for per-document accuracy or real-time inference, we optimize for \emph{corpus-scale batch annotation}: the pipeline should process tens of thousands of documents reliably, with minimal per-language setup cost. This leads to specific trade-offs -- prompt-based extraction instead of trained classifiers (no annotated data required), global deduplication instead of per-document processing (throughput over latency), and graph-enriched candidate generation instead of flat retrieval (recall over speed) -- that we believe are well-suited to the dataset-creation setting but would not necessarily be optimal for a real-time skill-matching API.

From a resource perspective, we operate entirely in Ukrainian using the official ESCO localization--an approach unexplored in the academic literature. Existing Ukrainian ESCO efforts remain proprietary, and the broader skill-extraction literature is overwhelmingly English-focused. By releasing both a raw corpus and an enriched version with document-level ESCO annotations, we provide the first openly available resource for Ukrainian labour-market skill analytics.

We note that our evaluation strategy differs from prior work: \citet{Clavi2023} and \citet{Decorte2025} evaluate on SkillSpan-ESCO with quantitative R-Precision and MRR metrics, while SkiLLMo uses domain-expert precision estimates. These benchmarks are English-only and span-level, making direct comparison infeasible for our Ukrainian, document-level setting. We instead adopt a multi-faceted evaluation (Section~\ref{sec:eval}) combining expert review, consistency analysis, and intrinsic metrics, designed to assess annotation quality at dataset scale rather than per-sentence extraction accuracy.

\section{Data Collection and Preprocessing}

To construct a large-scale, real-world dataset for studying the skill gap between labor market supply and demand, we collected data from **work.ua**, the largest Ukrainian online employment platform. work.ua hosts hundreds of thousands of active job postings and candidate resumes spanning a wide range of occupational categories, making it a representative source for analyzing skill requirements and offerings in the Ukrainian labor market. Data was collected over  June 2025 using a custom scraping pipeline built on the Scrapy framework with headless browser automation for dynamically rendered content.

The collection pipeline operates in two stages per entity type -- listing then detail -- applied independently to vacancies and resumes. In the listing stage, dedicated spiders paginate through work.ua search results, extracting structured ID lists for all available vacancies and resumes. In the detail stage, individual pages are fetched for each collected ID, yielding rich structured records. Each vacancy record contains approximately 20 fields spanning job title, salary and salary details, company metadata (name, size, industry), location, remote-work and employment-type indicators, an explicit skills list from the platform's structured skill tags, and the full job description as both HTML and plain text. Each resume record contains 13 fields: candidate title, name, age, city, desired salary, employment type, work-location preference, driver-license and own-car indicators, creation date, full profile content converted to Markdown, cross-links to the candidate's other resumes on the platform, and raw HTML. All text fields undergo whitespace normalization, and skill lists are filtered to remove empty entries. The resulting dataset comprises 98.4K vacancy records and 104K resume records after deduplication.

Both vacancy and resume documents contain two complementary skill representations: an explicit structured skills list curated by the poster, and free-text descriptions from which skills can be extracted using NLP methods. The cross-linking between a candidate's multiple resumes additionally enables analysis of how individuals present different skill sets across roles. This dual representation supports evaluation of both closed-vocabulary (taxonomy-based) and open-vocabulary (extraction-based) approaches to skill identification. The dataset was collected with explicit permission from work.ua and handled in full accordance with the platform's privacy policy. Resume records were anonymized by removing candidate names and contact details; the remaining fields reflect only information that candidates voluntarily publish on a public job-search platform. The dataset will be released upon publication under [LICENSE].

\begin{figure}[ht]
    \centering
    \includegraphics[width=0.9\linewidth]{dataset_work.png} 
    \caption{Raw dataset collection pipeline}
    \label{fig:dataseet1}
\end{figure}


\section{Pipeline Architecture}
\label{sec:pipeline}

Our pipeline processes documents in three sequential stages, illustrated in Figure~\ref{fig:pipeline}. All generative steps use Qwen3-30B-A3B \citep{yang2025qwen3technicalreport} -- a mixture-of-experts model with 30B total and 3B active parameters -- served locally via vLLM (model identifier: \texttt{Qwen/Qwen3-30B-A3B-Instruct-2507-FP8}). Embedding-based retrieval uses Gemini embeddings from \texttt{gemini-embedding-001} as this model is substantially outperforming prior embedding models on multilingual retrieval and classification tasks according to Massive Multilingual Text Embedding Benchmark (MMTEB) \citep{lee2025geminiembeddinggeneralizableembeddings}. We pre-compute and cache ESCO label and description embeddings offline, and compute query (skill-string) embeddings on demand with disk caching to avoid repeated API calls. All LLM prompts are in Ukrainian, matching both the input data and the ESCO localization; their full text is provided in Appendix~\ref{app:prompts}.

\begin{table}[t]
\centering
\small
\begin{tabular}{@{}ll@{}}
\toprule
\textbf{Parameter} & \textbf{Value} \\
\midrule
FUZZY\_THRESHOLD & 80 \\
EMBEDDING\_THRESHOLD & 0.72 \\
EMBEDDING\_TITLE\_WEIGHT & 0.6 \\
EMBEDDING\_SIBLING\_THRESHOLD & 0.5 \\
LLM\_MAX\_CANDIDATES & 100 \\
LLM\_BATCH\_SIZE & 10 \\
LLM\_MAX\_CONCURRENT & 32 \\
LLM\_RESCORE\_THRESHOLD & 0.7 \\
VLLM\_CANDIDATES\_PER\_SKILL & 10 \\
CV\_EXPERIENCE\_WEIGHT & 0.8 \\
CV\_EDUCATION\_WEIGHT & 0.5 \\
CV\_OVERLAP\_BOOST & 1.3 \\
\bottomrule
\end{tabular}
\caption{Key hyperparameters used in our experiments.}
\label{tab:hyperparams}
\end{table}

\begin{figure}[ht]
    \centering
    \includegraphics[width=0.9\linewidth]{pipeline_vrt.png} 
    \caption{Extraction pipeline architecture}
    \label{fig:pipeline}
\end{figure}

\subsection{Stage 1: Skill Extraction}
\label{sec:stage1}

For each document, the LLM extracts skill mentions from the available text fields, returning a structured JSON array of skill strings. Vacancies and resumes use separate extraction prompts tailored to their structure: the vacancy prompt operates over the job title and description, while the resume prompt receives labelled sections (desired position, work experiences with responsibilities, education, and additional information), allowing the model to extract skills from each section independently.

Extraction calls are fired concurrently across all documents in the batch (up to 32 parallel requests) and cached on disk by SHA-256 content hash, making re-runs on repeated inputs cost-free. When a document already carries platform-provided skill tags (\texttt{raw\_skills}), we include them verbatim as additional skill strings alongside the LLM-extracted skills.

\subsection{Stage 2: Retrieval-Augmented Global Mapping}
\label{sec:stage2}

After extraction, all unique skill strings across the document batch are collected via set union. The pipeline then maps each unique skill to an ESCO~URI through three sub-steps: multi-source candidate retrieval, graph enrichment, and LLM-based selection.

\paragraph{Multi-source candidate retrieval.}
For each unique skill, candidates are drawn from two complementary sources:

\begin{itemize}
\setlength\itemsep{0.2em}
\item \textbf{Fuzzy matching}: \texttt{rapidfuzz} token-sort-ratio is computed against all ESCO preferred labels (lowercased, punctuation-stripped). Matches scoring $\geq 80$ are retained, with confidence set to the normalized score. Near-exact matches (score $\geq 0.95$) receive a $+0.15$ boost to anchor precise lexical hits against generic embedding matches.

\item \textbf{Embedding similarity}: each skill is looked up against two pre-computed caches -- ESCO label embeddings and ESCO description embeddings (Gemini \texttt{gemini-embedding-001}, float32, L2-normalized). The final score per URI is a weighted combination: $s = 0.6 \times \text{label\_sim} + 0.4 \times \text{desc\_sim}$. The top-3 URIs per skill with $s \geq 0.72$ become embedding candidates.
\end{itemize}

\paragraph{ESCO graph enrichment.}
For each candidate URI retrieved by either method, the system traverses the ESCO directed acyclic graph. \emph{Parent} nodes are added with the child's retrieval score. \emph{Sibling} nodes (children of the same parent) are added only if their embedding similarity to the source skill exceeds $0.5$, scored at the computed similarity. This enrichment step expands the candidate pool with semantically adjacent concepts that neither fuzzy nor embedding retrieval would surface directly -- for example, retrieving ``manage human resources'' when the raw mention is ``HR management'' and only ``plan human resource management'' was directly retrieved.

\paragraph{Category-aware candidate selection.}
To prevent any single retrieval source from dominating the candidate pool, the selection algorithm guarantees minimum slots per category: 15 fuzzy, 15 embedding, 5 graph-sibling, and 3 graph-parent candidates. Remaining slots (up to 100 total) are filled from a globally score-sorted pool. The final candidate list is serialized as numbered lines of the form \texttt{- [idx] Label: ... | Alt: ...}, with ESCO~URIs deliberately omitted. The LLM returns an integer index, which is resolved to a URI via a lookup dictionary -- making it structurally impossible for the model to fabricate URIs.

\paragraph{Batch LLM mapping.}
Skills are grouped into batches of 10 and sent to the LLM with the mapping prompt and the shared candidate list. All batches are dispatched concurrently (up to 32 parallel requests) and cached by content hash. The LLM returns a JSON structure containing mappings (with skill, candidate index, label, confidence, and reasoning) and a list of unmapped skills. Every returned index is validated against \texttt{uri\_exists()} before acceptance; invalid indices are discarded silently.

\subsection{Stage 3: Context-Aware Rescoring}
\label{sec:stage3}

Global mapping (Stage~2) trades document context for efficiency: a skill is mapped once regardless of where it appears. This creates ambiguity for polysemous or context-dependent terms. Stage~3 recovers this lost context selectively.

For each document, the pipeline identifies skills whose Stage~2 confidence falls below $0.7$ or that remained unmapped. If such skills exist and the document carries title/description text, a rescoring prompt is dispatched containing the document title, a truncated description ($\leq$1000 characters), the skills to rescore, and the same candidate list used in Stage~2. The LLM re-evaluates each skill in the specific document context -- for example, disambiguating ``Java'' as a programming language rather than a geographic reference when the vacancy title is ``Backend Developer.''

Rescoring calls are fired concurrently across all documents requiring them. A rescored mapping replaces the Stage~2 result \emph{only if it achieves strictly higher confidence}, ensuring that rescoring never degrades a mapping that Stage~2 already handled well.

\subsection{Baselines}
\label{sec:baselines}

We compare the three-stage pipeline against simpler approaches implemented within the same pipeline scaffold:

\begin{itemize}
\setlength\itemsep{0.2em}
\item \textbf{FuzzyMapper}: standalone \texttt{rapidfuzz} matching against ESCO labels. No LLM, no embeddings, no graph enrichment.
\item \textbf{EmbeddingMapper}: standalone embedding cosine similarity with the weighted label+description score. No LLM re-ranking.
\item \textbf{LLM direct} (\texttt{llm\_direct}): single-pass LLM mapping with the same candidate retrieval and graph enrichment as the full pipeline, but processing each document independently without cross-document deduplication or a separate rescoring stage.
\item \textbf{vLLM optimized} (\texttt{vllm\_optimized}): per-skill LLM mapping optimized for throughput by sending one skill with a small focused candidate list (10 candidates: fuzzy top-5 + embedding top-5), using a constant system message to maximize vLLM prefix-cache reuse; generation is capped at 512 tokens per request.
\item \textbf{CV weighted} (\texttt{cv\_weighted}): a resume-only multi-signal mapper that retrieves ESCO candidates from structured CV sections via embedding search (work experience position/responsibilities top-10 with weight 0.8 and overlap boost 1.3; education faculty top-5 with weight 0.5), merging these with platform skills.
\end{itemize}

\section{Evaluation}
\label{sec:eval}
Evaluating taxonomy normalization pipelines against massive ontologies like ESCO presents a known methodological challenge. The taxonomy contains nearly 14,000 deeply granular and frequently overlapping nodes (e.g., distinguishing between "Python (programming language)" and "scripting in Python"). Consequently, constructing a manual human-annotated "golden set" is not only prohibitively expensive but highly subjective, often suffering from low inter-annotator agreement.

To rigorously validate our pipeline without relying on intractable manual annotation, we implement a multi-pronged evaluation strategy: (1)~expert judgment on a stratified sample of mappings, (2)~calibration of an LLM-as-judge against these expert annotations, (3)~intrinsic metrics at full scale, (4)~baseline comparison of mapper strategies, (5)~ablation of pipeline stages, and (6)~mapping consistency analysis.

\subsection{Expert Evaluation}
\label{sec:eval-expert}

To assess mapping quality, we construct a stratified evaluation dataset of 447 unique skill mappings drawn from the full pipeline output (8,176 vacancies and 3,984 resumes processed with \texttt{llm\_two\_stage}). The 447 mappings are stratified by corpus frequency into three bands of 149 skills each: \emph{frequent} skills (the most common raw skill strings by occurrence count, representing the mappings with the highest aggregate impact on dataset quality), \emph{mid-frequency} skills (randomly sampled from skills appearing 30--100 times, representing the middle of the distribution), and \emph{rare} skills (randomly sampled from skills appearing exactly once, typically long domain-specific phrases extracted by the LLM from job descriptions). This stratification ensures that the evaluation covers both high-impact generic skills (e.g., \uk{відповідальність} $\to$ \uk{нести відповідальність}) and challenging tail cases (e.g., \uk{монтаж каналізації зовнішньої} $\to$ \uk{прокладати труби}) where the pipeline is most likely to err.

Each mapping is presented to three independent domain experts as a (raw skill string, ESCO label) pair. Experts assign one of three judgments: \emph{Correct} (exact or near-exact semantic match), \emph{Acceptable} (related but imprecise -- e.g., mapped to a parent concept or close sibling), or \emph{Incorrect} (wrong or unrelated ESCO skill). This three-point ordinal scale is standard in taxonomy annotation tasks where binary correct/incorrect judgments are too coarse to capture near-miss mappings to parent or sibling concepts \citep{Malandri2025}.

From these judgments we derive two accuracy metrics: \emph{strict accuracy} (fraction rated Correct) and \emph{lenient accuracy} (fraction rated Correct or Acceptable), reported per frequency band and overall. For aggregation across the three experts, we use majority vote when a majority exists; otherwise (three-way disagreement) we use an ordinal-median tie-break. Inter-annotator agreement is measured using Fleiss' $\kappa$.

\paragraph{Results.} Table~\ref{tab:expert-eval} reports expert evaluation results across the three frequency bands.

\begin{table}[h]
\centering
\small
\begin{tabular}{@{}lrccc@{}}
\toprule
\textbf{Band} & \textbf{N} & \textbf{Strict} & \textbf{Lenient} & \textbf{$\kappa$} \\
\midrule
Frequent & 149 & 82.6\% & 96.6\% & 0.371 \\
Mid-frequency & 149 & 81.2\% & 96.0\% & 0.387 \\
Rare & 149 & 65.8\% & 91.9\% & 0.232 \\
\midrule
\textbf{Overall} & \textbf{447} & \textbf{76.5\%} & \textbf{94.9\%} & \textbf{0.332} \\
\bottomrule
\end{tabular}
\caption{Expert evaluation of mapping quality by frequency band. Strict = fraction rated Correct; Lenient = Correct + Acceptable. $\kappa$ = Fleiss' kappa (3 experts, majority-vote aggregation with ordinal-median tie-break).}
\label{tab:expert-eval}
\end{table}

The pipeline achieves 94.9\% lenient accuracy overall, indicating that the vast majority of mappings are judged as semantically correct or acceptably close by domain experts. Strict accuracy is 76.5\%, with a clear gradient across frequency bands: frequent skills (82.6\%) and mid-frequency skills (81.2\%) are mapped with high precision, while rare skills (65.8\%) prove more challenging. This is expected -- frequent skills tend to have direct ESCO label matches, while rare skills are often long, domain-specific LLM-extracted phrases (e.g., \uk{забезпечення точності при роботі з документами} $\to$ \uk{перевіряти правильність інформації}) that map to a semantically related but not identical ESCO concept. Notably, lenient accuracy remains above 90\% even for rare skills, confirming that the pipeline consistently identifies the correct semantic neighborhood even when exact label matching fails.

Inter-annotator agreement (Fleiss' $\kappa$ = 0.332) reflects the inherent subjectivity of judging semantic equivalence across a 14,000-node taxonomy, particularly at the Correct/Acceptable boundary where reasonable experts may disagree on whether a mapping to a parent concept or close sibling constitutes an exact match or merely an acceptable one. Raw agreement ranges from 70\% (frequent/mid) to 43\% (rare), with the drop driven by the greater ambiguity of long domain-specific phrases rather than annotator unreliability.
\paragraph{Limitations.} A limitation of this evaluation design is that experts judge mappings \emph{without the original document context}. Stage~3 rescoring can produce mappings that are justified by a specific vacancy but appear incorrect in isolation -- for example, \uk{проектування нових виробів} (new product design) mapped to \uk{3D-моделювання} (3D modeling) seems wrong out of context, but is defensible for an industrial design vacancy where new product development is carried out primarily through 3D modeling. Such context-dependent mappings may be penalized by evaluators. The evaluation therefore provides a conservative estimate of mapping quality; the true accuracy of the full pipeline with Stage~3 rescoring is likely higher for context-sensitive skills.

\subsection{LLM-as-Judge Validation}
\label{sec:eval-judge}

To enable scalable evaluation across all mapper configurations, we calibrate an LLM judge (Gemini~2.5~Flash) against the 447 expert-annotated mappings from Section~\ref{sec:eval-expert}. The judge receives each (raw skill, ESCO label) pair and assigns the same three-point ordinal rating (Correct / Acceptable / Incorrect).

On lenient accuracy the judge agrees with the expert majority vote 86.4\% of the time; strict agreement is 55.0\% (Cohen's $\kappa = 0.237$). Inspection of the confusion matrix reveals a systematic conservative bias: 132 mappings rated Correct by experts are downgraded to Acceptable by the judge, while only 8 are upgraded in the opposite direction. The judge is therefore a pessimistic proxy -- it underestimates absolute quality but is expected to preserve the \emph{relative} ranking of mappers, since the conservative bias applies uniformly across mapping strategies. We use it in Section~\ref{sec:eval-baseline} to compare mapper strategies on a 50-document sample where expert annotation would be prohibitively expensive.

\subsection{Intrinsic Metrics}
\label{sec:eval-intrinsic}

Table~\ref{tab:intrinsic} reports aggregate statistics for the full production pipeline (\texttt{llm\_two\_stage}) applied to 8,176 vacancies and 3,984 resumes. These numbers characterize the released dataset; controlled comparisons across mapper strategies and pipeline stages follow in Sections~\ref{sec:eval-baseline}--\ref{sec:eval-ablation}.

\begin{table}[h]
  \centering
  \small
  \begin{tabular}{@{}lrr@{}}
    \toprule
    \textbf{Metric} & \textbf{Vacancies} & \textbf{Resumes} \\
    \midrule
    Documents & 8{,}176 & 3{,}984 \\
    Unmapped rate & 0.02\% & 0.13\% \\
    Mean confidence & 0.877 & 0.879 \\
    Median confidence & 0.900 & 0.900 \\
    Skills per document & 29.0 & 48.9 \\
    Distinct ESCO URIs & 4{,}877 & 2{,}916 \\
    \bottomrule
  \end{tabular}
  \caption{Production-scale dataset statistics for the full pipeline on vacancies and resumes.}
  \label{tab:intrinsic}
\end{table}

The pipeline achieves near-zero unmapped rates (0.02\% for vacancies, 0.13\% for resumes) with mean confidence scores of 0.877--0.879. Vacancies produce 29.0 skills per document and cover 4{,}877 distinct ESCO URIs; resumes yield 48.9 skills per document but map to 2{,}916 URIs. Table~\ref{tab:confidence} shows that approximately 90\% of all mappings receive confidence $\geq 0.8$ for both data sources.

\begin{table}[h]
  \centering
  \small
  \begin{tabular}{@{}lrr@{}}
    \toprule
    \textbf{Confidence range} & \textbf{Vacancies} & \textbf{Resumes} \\
    \midrule
    0.9--1.0 & 57.5\% & 60.1\% \\
    0.8--0.9 & 32.0\% & 29.0\% \\
    0.7--0.8 & 10.0\% & 9.4\% \\
    0.6--0.7 & 0.4\% & 1.1\% \\
    $<$0.6 & 0.1\% & 0.4\% \\
    \bottomrule
  \end{tabular}
  \caption{Distribution of mapping confidence scores.}
  \label{tab:confidence}
\end{table}

\subsection{Baseline Comparison}
\label{sec:eval-baseline}

We compare four mapper strategies on a sample of 50 vacancies and 45 resumes, reporting both intrinsic metrics and LLM judge accuracy (calibrated in Section~\ref{sec:eval-judge}). We introduce \emph{effective correct mappings per document} (mappings/doc $\times$ strict accuracy) as a composite metric capturing both coverage and precision.

An important caveat: the mappers differ not only in their mapping strategy but also in the extraction stage they use. Fuzzy and LLM direct operate on platform-provided skills only (4.2 skills/doc for vacancies, 10.1 for resumes), while Embedding and LLM two-stage additionally incorporate LLM-extracted skills from document text ($\sim$20 skills/doc for vacancies, $\sim$34 for resumes). This reflects the full pipeline design -- fuzzy and LLM direct serve as mapping-only baselines, while Embedding and LLM two-stage are evaluated as end-to-end systems. The ``Skills'' column in Tables~\ref{tab:baseline-vac}--\ref{tab:baseline-res} makes this difference explicit.

\begin{table}[h]
  \centering
  \small
  \begin{tabular}{@{}lrrrrrr@{}}
    \toprule
    \textbf{Mapper} & \textbf{Skills} & \textbf{Unmap.} & \textbf{Conf.} & \textbf{Strict} & \textbf{Lenient} & \textbf{Eff./doc} \\
    \midrule
    Fuzzy & 4.2 & 75.9\% & 0.857 & 86.1\% & 90.7\% & 0.74 \\
    Embedding & 20.8 & 4.1\% & 0.806 & 51.0\% & 79.0\% & 10.15 \\
    LLM direct & 4.2 & 10.3\% & 0.931 & 90.0\% & 99.0\% & 3.29 \\
    LLM two-stage & 19.8 & 29.1\% & 0.905 & 69.0\% & 94.0\% & 9.16 \\
    \bottomrule
  \end{tabular}
  \caption{Baseline mapper comparison on vacancies (50-doc sample). Skills = mean input skills/doc; Strict/Lenient = LLM judge accuracy; Eff./doc = effective correct mappings per document.}
  \label{tab:baseline-vac}
\end{table}

\begin{table}[h]
  \centering
  \small
  \begin{tabular}{@{}lrrrrrr@{}}
    \toprule
    \textbf{Mapper} & \textbf{Skills} & \textbf{Unmap.} & \textbf{Conf.} & \textbf{Strict} & \textbf{Lenient} & \textbf{Eff./doc} \\
    \midrule
    Fuzzy & 10.1 & 90.0\% & 0.913 & 69.4\% & 85.7\% & 0.76 \\
    Embedding & 34.9 & 6.4\% & 0.808 & 64.0\% & 89.0\% & 20.34 \\
    LLM direct & 10.1 & 27.0\% & 0.924 & 86.0\% & 97.0\% & 6.10 \\
    LLM two-stage & 33.9 & 29.6\% & 0.893 & 54.0\% & 95.0\% & 13.33 \\
    \bottomrule
  \end{tabular}
  \caption{Baseline mapper comparison on resumes (45-doc sample). Columns as in Table~\ref{tab:baseline-vac}.}
  \label{tab:baseline-res}
\end{table}

The results reveal clear trade-offs. Among the mapping-only baselines (platform skills only), \textbf{fuzzy matching} achieves high precision on the few skills it maps (86.1\%/69.4\% strict) but leaves 76--90\% of skills unmapped, producing under 1 effective correct mapping per document. \textbf{LLM direct} yields the highest per-mapping precision (86--90\% strict, 97--99\% lenient) but maps only 3.7--7.1 skills per document due to its single-pass per-document design. Among the end-to-end systems (with LLM extraction), \textbf{embedding search} achieves near-complete coverage (4--6\% unmapped) but sacrifices precision (51--64\% strict), as cosine similarity alone cannot distinguish semantically close but taxonomically distinct ESCO nodes. The \textbf{LLM two-stage} pipeline achieves the highest lenient accuracy (94--95\%) among end-to-end systems with 9.16--13.33 effective correct mappings per document, combining the broad candidate retrieval of embedding search with LLM re-ranking to filter false positives. On resumes, embedding search produces more effective correct mappings per document (20.34 vs.\ 13.33) owing to its lower unmapped rate, but at the cost of substantially lower lenient accuracy (89\% vs.\ 95\%).

\subsection{Ablation Study}
\label{sec:eval-ablation}

To quantify the contribution of each pipeline stage, we run four configurations on a 50-document sample: Stage~1 only (extraction), Stage~1+2 without graph enrichment, Stage~1+2 with graph enrichment, and the full three-stage pipeline.

\begin{table}[h]
  \centering
  \small
  \begin{tabular}{@{}lrrrrr@{}}
    \toprule
    \textbf{Configuration} & \textbf{Skills} & \textbf{Unmap.} & \textbf{Conf.} & \textbf{Graph} & \textbf{Cons.} \\
    \midrule
    Stage 1 only & 19.8 & --- & --- & --- & --- \\
    Stage 1+2 (no graph) & 19.8 & 37.7\% & 0.908 & 0.0\% & 1.000 \\
    Stage 1+2 (w/ graph) & 19.8 & 44.4\% & 0.918 & 23.9\% & 1.000 \\
    Full pipeline & 19.8 & 31.7\% & 0.899 & 24.3\% & 0.892 \\
    \bottomrule
  \end{tabular}
  \caption{Ablation study on vacancies (50-doc sample). Skills = mean skills/doc; Unmap.\ = unmapped rate; Conf.\ = mean confidence; Graph = graph enrichment yield; Cons.\ = mapping consistency.}
  \label{tab:ablation-vac}
\end{table}

\begin{table}[h]
  \centering
  \small
  \begin{tabular}{@{}lrrrrr@{}}
    \toprule
    \textbf{Configuration} & \textbf{Skills} & \textbf{Unmap.} & \textbf{Conf.} & \textbf{Graph} & \textbf{Cons.} \\
    \midrule
    Stage 1 only & 33.9 & --- & --- & --- & --- \\
    Stage 1+2 (no graph) & 33.9 & 47.7\% & 0.909 & 0.0\% & 1.000 \\
    Stage 1+2 (w/ graph) & 33.9 & 41.3\% & 0.906 & 15.2\% & 1.000 \\
    Full pipeline & 33.9 & 30.1\% & 0.898 & 16.0\% & 0.919 \\
    \bottomrule
  \end{tabular}
  \caption{Ablation study on resumes (50-doc sample). Columns as in Table~\ref{tab:ablation-vac}.}
  \label{tab:ablation-res}
\end{table}

The results show two distinct effects. \emph{Graph enrichment} reduces unmapped skills for resumes (47.7\% $\to$ 41.3\%) by expanding candidates via ESCO parent and sibling relations -- consistent with the 20.4\% graph enrichment yield observed at full scale (Table~\ref{tab:intrinsic}). For vacancies, however, graph enrichment \emph{increases} the unmapped rate (37.7\% $\to$ 44.4\%): the expanded candidate list makes the LLM's ranking task harder, and for explicitly stated vacancy skills that already have strong direct matches, the additional graph candidates introduce noise rather than useful alternatives.

\emph{Stage~3 rescoring} then substantially reduces unmapped rates in both cases: from 44.4\% to 31.7\% on vacancies and from 41.3\% to 30.1\% on resumes. The combined effect of graph enrichment and Stage~3 reduces the overall unmapped rate from 37.7\% to 31.7\% (vacancies) and from 47.7\% to 30.1\% (resumes).

The trade-off is a modest decrease in mean confidence (0.908 $\to$ 0.899 for vacancies; 0.909 $\to$ 0.898 for resumes) and consistency (1.000 $\to$ 0.892/0.919). Both effects are expected: Stage~3 introduces document-level context that can shift a skill's mapping depending on the surrounding vacancy or resume, producing legitimate context-dependent variation. Without Stage~3, the pipeline maps each skill identically regardless of context (consistency = 1.000), which maximizes consistency at the cost of ignoring contextual disambiguation.

\subsection{Consistency Analysis}
\label{sec:eval-consistency}

We measure mapping consistency as the fraction of occurrences of a skill that map to the modal ESCO URI, computed over skills appearing at least 5 times in the corpus.

\begin{table}[h]
  \centering
  \small
  \begin{tabular}{@{}lrr@{}}
    \toprule
    \textbf{Metric} & \textbf{Vacancies} & \textbf{Resumes} \\
    \midrule
    Total unique skills & 86{,}093 & 111{,}027 \\
    Qualifying skills ($\geq$5 occ.) & 4{,}769 & 3{,}455 \\
    Macro-avg consistency & 0.330 & 0.443 \\
    \midrule
    Polysemy count & 1{,}731 & 1{,}237 \\
    Instability count & 3{,}002 & 1{,}883 \\
    \midrule
    Consistency = 100\% & 36 & 335 \\
    Consistency 90--99\% & 0 & 58 \\
    Consistency 80--89\% & 63 & 164 \\
    Consistency 70--79\% & 49 & 110 \\
    Consistency 60--69\% & 263 & 232 \\
    Consistency 50--59\% & 304 & 266 \\
    Consistency $<$50\% & 4{,}054 & 2{,}290 \\
    \bottomrule
  \end{tabular}
  \caption{Mapping consistency analysis. Polysemy = skills legitimately mapping to multiple ESCO concepts; Instability = skills with inconsistent mappings not attributable to polysemy.}
  \label{tab:consistency}
\end{table}

Macro-average consistency is 0.330 for vacancies and 0.443 for resumes. While these numbers appear low, they reflect two distinct phenomena. First, \emph{genuine polysemy}: many Ukrainian skill terms are ambiguous and legitimately map to different ESCO nodes depending on context. For instance, \uk{м.е.док} (M.E.Doc, Ukrainian accounting software) appears 32 times and maps to 14 different ESCO URIs spanning digital skills, document management, and accounting -- all defensible depending on the vacancy. Second, \emph{context-dependent variation}: Stage~3 rescoring adapts mappings to each document's context, so the same skill string maps differently depending on the surrounding vacancy or resume -- as confirmed by the ablation study (Section~\ref{sec:eval-ablation}), where removing Stage~3 yields perfect consistency (1.000).

Resumes exhibit higher consistency (335 skills at 100\% vs.\ 36 for vacancies), reflecting more standardized skill descriptions in CVs. The consistency--coverage trade-off mirrors the ablation findings: Stage~3 rescoring sacrifices some consistency to reduce unmapped rates, a trade-off we consider favorable given that context-aware disambiguation is a design goal of the pipeline.

\section*{Limitations}



\section*{Acknowledgments}



% Bibliography entries for the entire Anthology, followed by custom entries
%\bibliography{anthology,custom}
% Custom bibliography entries only
\bibliography{custom}

\appendix

\section{Pipeline Prompts}
\label{app:prompts}

All prompts used in the pipeline are provided below in their original Ukrainian form. Template variables (e.g., \texttt{\{title\}}, \texttt{\{skills\}}, \texttt{\{candidates\}}) are filled at runtime. Baseline prompts are available in the project repository.

\subsection{Stage~1 --- Skill Extraction}

\paragraph{Vacancy Extraction Prompt.}
Used to extract skills from vacancy title and description text.

\begin{small}
\begin{alltt}\fontencoding{T2A}\selectfont
Ви — експерт з аналізу ринку праці та професійних навичок.

Витягніть усі навички, компетентності та кваліфікації,
згадані в тексті вакансії нижче.
Включайте як явні навички (наприклад, "Python",
"управління проєктами"), так і неявні (наприклад,
"вміння працювати в команді").

Назва вакансії: \{title\}
Текст вакансії: \{text\}

Поверніть JSON-об'єкт:
\{
  "skills": [
    "<навичка 1>",
    "<навичка 2>",
    ...
  ]
\}

Витягуйте лише реальні навички та компетентності
українською мовою. Не включайте вимоги до досвіду
на кшталт "3 роки досвіду", якщо вони не вказують
на конкретну навичку.
\end{alltt}
\end{small}

\paragraph{CV Extraction Prompt.}
Used to extract skills from structured resume sections.

\begin{small}
\begin{alltt}\fontencoding{T2A}\selectfont\itshape
Ви — експерт з аналізу ринку праці та професійних навичок.

Витягніть усі навички, компетентності та кваліфікації,
згадані у резюме нижче.
Включайте як явні навички (наприклад, "Python",
"управління проєктами"), так і неявні (наприклад,
"вміння працювати в команді").

Бажана посада: \{title\}

Досвід роботи:
\{work\_experiences\}

Освіта:
\{educations\}

Додаткова інформація:
\{additional\_info\}

Поверніть JSON-об'єкт:
\{
  "skills": [
    "<навичка 1>",
    "<навичка 2>",
    ...
  ]
\}

Витягуйте лише реальні навички та компетентності
українською мовою. Не включайте вимоги до досвіду
на кшталт "3 роки досвіду", якщо вони не вказують
на конкретну навичку.
\end{alltt}
\end{small}

\subsection{Stage~2 --- Global Mapping}

\paragraph{Mapping Prompt.}
Maps a batch of extracted skills to numbered ESCO candidates. Used by the two-stage mapper.

\begin{small}
\begin{alltt}\fontencoding{T2A}\selectfont\itshape
Ви — експерт з професійних навичок та таксономії ESCO.

Зіставте кожну витягнуту навичку з найбільш відповідною
навичкою ESCO серед запропонованих кандидатів.

Кандидати пронумеровані (наприклад, [1], [2], ...).
Поверніть номер кандидата у полі "esco\_id".

Витягнуті навички: \{skills\}

Кандидати ESCO:
\{candidates\}

Поверніть JSON-об'єкт:
\{
  "mappings": [
    \{
      "raw\_skill": "<витягнута навичка>",
      "esco\_id": <номер кандидата>,
      "esco\_label": "<бажана мітка ESCO>",
      "confidence": <0.0-1.0>,
      "reasoning": "<коротке пояснення>"
    \}
  ],
  "unmapped": ["<навички, для яких не знайдено
                 впевненого відповідника>"]
\}

Зіставляйте лише за наявності чіткого семантичного
збігу. Значення confidence має відображати вашу
впевненість.
\end{alltt}
\end{small}

\subsection{Stage~3 --- Context-Aware Rescoring}

\paragraph{Rescore with Context Prompt.}
Re-scores ambiguous mappings using the full vacancy context (title and description).

\begin{small}
\begin{alltt}\fontencoding{T2A}\selectfont\itshape
Ви — експерт з професійних навичок та таксономії ESCO.

Контекст вакансії:
Назва: \{title\}
Опис: \{description\}

Повторно оцініть зіставлення навичок з ESCO,
враховуючи контекст вакансії.
Для кожної навички виберіть найкращого кандидата ESCO
або позначте як незіставлену.

Кандидати пронумеровані (наприклад, [1], [2], ...).
Поверніть номер кандидата у полі "esco\_id".

Навички для переоцінки: \{skills\}

Кандидати ESCO:
\{candidates\}

Поверніть JSON-об'єкт:
\{
  "mappings": [
    \{
      "raw\_skill": "<навичка>",
      "esco\_id": <номер кандидата>,
      "esco\_label": "<бажана мітка ESCO>",
      "confidence": <0.0-1.0>,
      "reasoning": "<коротке пояснення з урахуванням
                     контексту вакансії>"
    \}
  ],
  "unmapped": ["<навички, для яких не знайдено
                 відповідника>"]
\}

Зіставляйте лише за наявності чіткого семантичного
збігу. Значення confidence має відображати вашу
впевненість.
\end{alltt}
\end{small}

\end{document}
